\chapter{Research and Existing Solutions}

\section{Research Method / Approach}

The research combined reading with measurement on our own data, and each stage fed a design decision (Section~2.8).

\begin{enumerate}
\item \textbf{Literature and software review.} We reviewed maintenance-management software categories, keyword and semantic retrieval, retrieval-augmented generation \citep{asai2023,lewis2020}, graph-based retrieval \citep{edge2024} and case-based reasoning \citep{aamodt1994}. The review is architectural: it compares how approaches are designed and does not claim measured superiority of any of them, nor does it make claims about specific commercial products.
\item \textbf{Data profiling.} We downloaded the 2025 FAA SDR file and profiled it with \texttt{scripts/fetch\_sdr\_data.py}: 67{,}637 raw rows, 66{,}658 after removing duplicates (1.45\%), spread over 8{,}550 aircraft with a median of 4 and a maximum of 153 records per aircraft.
\item \textbf{Ground truth from the data itself.} Reporters sometimes file a supplemental report that begins ``SUPPLEMENTAL REPORT FOR (record id)''. Extracting these gave 141 valid correction pairs (six more were discarded: five cite a record outside the corpus, one crosses aircraft). No labels were written by the team \citep{faa2025}.
\item \textbf{A working prototype and measurement.} We built the pipeline, measured the conflict classifier on the 141 positive and 209 negative pairs, and measured the cost of running it over the corpus (Section~2.7).
\end{enumerate}

\section{Existing Solutions / Products}

\subsection{Maintenance-management systems}

Computerised maintenance management systems administer assets, work orders, schedules and inspection history. Their strength is record keeping and workflow. Searching history is usually by field or keyword, and a later correction is a separate record that the system does not connect to the one it corrects. We do not review specific vendors; a formal product comparison is outside this iteration.

\subsection{Keyword retrieval}

Matching literal terms is cheap and easy to explain, but two technicians describing the same condition in different words are not matched, and nothing in the result says whether a record was later corrected.

\subsection{Semantic retrieval and retrieval-augmented generation}

Embedding-based retrieval finds records whose meaning is close even when the words differ, and retrieval-augmented generation lets a language model answer from retrieved text \citep{lewis2020}. Variants add self-critique of the retrieved passages \citep{asai2023}. These methods rank by similarity; similarity says nothing about whether the retrieved intervention worked or has since been superseded.

\subsection{Graph-based retrieval}

Graph retrieval represents entities and their relationships explicitly and uses them to narrow and structure retrieval \citep{edge2024}. This is the closest published approach to ours: Halite adapts the idea to maintenance records, where the relationship that matters most is correction.

\subsection{Case-based reasoning}

Case-based reasoning solves a new problem by retrieving similar past cases and adapting their solutions \citep{aamodt1994}. Our diagnostic ranking is case-based: a symptom is compared with earlier records and the components they implicate are ranked. We add record validity (a superseded case counts for less) and evidence provenance (every case is a citable record).

\section{Comparison of Existing Solutions}

Table~\ref{tab:retrieval-compare} compares the approaches on criteria that matter for this problem. The last column states what is implemented in this iteration, not what is planned.

\begin{table}[H]
\centering
\caption{Architectural comparison of retrieval approaches}
\label{tab:retrieval-compare}
\footnotesize
\begin{tabularx}{\textwidth}{@{}>{\raggedright\arraybackslash}p{0.16\textwidth} >{\raggedright\arraybackslash}p{0.12\textwidth} >{\raggedright\arraybackslash}p{0.14\textwidth} >{\raggedright\arraybackslash}p{0.14\textwidth} L@{}}
\toprule
Capability & Keyword & Semantic RAG & Graph RAG & Halite (at Mid) \\
\midrule
Natural-language input & Limited & Yes & Yes & Yes (lexical matching only) \\
Asset-specific context & Manual & Filter & Graph & Graph (aircraft history plus capped fleet scan) \\
Historical corrections & No & Not intrinsic & Possible & Explicit SUPERSEDES edges; superseded records down-weighted \\
Ranked hypotheses & No & Generated text & Generated text & Deterministic, documented heuristic \\
Evidence checking & Manual & Partial & Partial & Every cited record identifier checked against the graph \\
Uncertainty handling & No & Model-based & Model-based & Heuristic labelled not-a-probability; unanalysed state shown; model failure reported \\
Semantic similarity & No & Yes & Yes & Not implemented (planned) \\
\bottomrule
\end{tabularx}
\end{table}


\section{Limitations and Gaps}

Existing approaches treat a record's relevance as similarity and its validity as someone else's problem. Three gaps follow.

\begin{enumerate}
\item \textbf{Validity is invisible.} A retrieved record may have been corrected by a later one; nothing tells the technician.
\item \textbf{Corrections change the component.} In the 141 FAA-recorded pairs, the superseding record names a different part from the one it supersedes in 92 cases (65.2\%), a different ATA/JASC code in 60 (42.6\%) and a different location in 67 (47.5\%). This does not show that the earlier report was wrong; it shows that the follow-up often points at a different component, so a system that ignores later records will often surface the wrong one.
\item \textbf{Generated text is not evidence.} A language model can cite a record that does not exist. Each cited identifier needs checking against the data.
\end{enumerate}

\section{Proposed Solution}

Halite records order and correction in a graph, retrieves evidence only from a bounded neighbourhood of that graph, ranks components with a deterministic heuristic, lets a model phrase (not decide) the answer, and checks every cited record identifier against the graph. Records that cannot be resolved with confidence go to a human reviewer. The design is described in Chapter~5.

\section{Complex Computing Problem}

How can a diagnostic system identify, rank and explain plausible fault hypotheses from heterogeneous, temporally evolving and possibly contradictory maintenance records, while preserving evidence provenance and avoiding unsupported high-confidence conclusions?

It is a complex problem for the following reasons, each of which we met in this iteration.

\begin{itemize}
\item \textbf{Heterogeneous text.} Free-text reports with inconsistent abbreviations, and fields such as component and location that are sometimes blank or UNKNOWN.
\item \textbf{Conflicting and evolving evidence.} A diagnosis can change after a later inspection, so order matters and a record's validity depends on records that came after it.
\item \textbf{A large candidate space.} Comparing every record with every earlier record is quadratic; the largest aircraft alone has 153 records and 11{,}628 pairs.
\item \textbf{Conflicting objectives.} Finding corrections (recall) trades against asserting false corrections (the baseline we measured asserts one on every pair).
\item \textbf{Evidence accountability.} Each displayed conclusion needs a traceable record.
\end{itemize}

\section{Technical Challenges}

These are challenges we encountered and what happened.

\begin{table}[H]
\centering
\caption{Technical challenges met in this iteration}
\label{tab:tech-challenges}
\small
\begin{tabularx}{\textwidth}{@{}>{\raggedright\arraybackslash}p{0.22\textwidth} L@{}}
\toprule
Challenge & What we found and did \\
\midrule
Query fan-out & The first service-record view issued one request per aircraft (1+118 on load). It now loads two page-initiated calls; the assets list is one aggregation query ($\approx$1.0\,s on staging). \\
Cost of conflict analysis & On the real corpus, 1{,}835 of 2{,}079 records have a candidate. Five-vote classification projects to 9{,}175 model requests vs a free-tier limit of 1{,}000/day. An opt-in cascade projects to 2{,}062--5{,}155 requests. Full-corpus token cost is not yet measured. \\
Model instability & On one pair, two models gave different labels. We record method, model and votes for every classification and do not tune thresholds to look better. \\
Legacy placeholder data & Early ingest carried confidence 1.0 without analysis. The API now reports such untouched values as unanalysed. \\
Provider limits & HTTP 429 is normal (27 on a 13-record run). Retries are bounded; a daily budget stops runs cleanly; a missing live model fails closed. \\
Empty model output & Staging fresh diagnoses returned no explanation text. Empty model output is an explicit HTTP 503, never a blank success. \\
Dense graphs & The largest aircraft produces a 305-edge subgraph; the interface shows one aircraft at a time. \\
\bottomrule
\end{tabularx}
\end{table}


\section{Research Findings and Design Decisions}

Table~\ref{tab:findings-design} maps each research finding from this iteration to the design decision it motivated.

\begin{table}[H]
\centering
\caption{How findings shaped the design}
\label{tab:findings-design}
\small
\begin{tabularx}{\textwidth}{@{}L L@{}}
\toprule
Finding & Design decision \\
\midrule
Corrections change the component in 65.2\% of FAA-recorded pairs & Represent SUPERSEDES explicitly; down-weight superseded records in ranking (factor 0.3); show the corrector next to a corrected record \\
Exhaustive pairing is quadratic & Only compare within the same aircraft, only with earlier records, only when at least one matching signal exists, at most 25 candidates \\
Models disagree and fail & Ranking is deterministic and needs no model; the model only phrases it; failures are reported, not hidden \\
Model-generated identifiers can be invented & Every identifier in the text is checked against the graph; unmatched ones are marked unverified \\
Analysis is expensive and rate-limited & Ingest stores records as unanalysed; analysis is a separate, resumable, budgeted step; the interface shows ``Not analysed'' honestly \\
Labels come free from the data & Use FAA-recorded supplemental pairs as ground truth rather than hand labels \\
\bottomrule
\end{tabularx}
\end{table}

